Throughout, $e$ is the physical-vector map $\hb^{-1}$ in components, $e^a(v) = \gamma^a \cdot \hb^{-1}(v)$, and $\omega$ the rotation-gauge field in the frame $\{\gamma_a\}$ as Cartan one-forms, whose values on vectors are the negative of the definition's $\Omega$ (dictionary convention C10; the ranks and the quadratic identity of (d) do not depend on that sign), so the component computation is one of Gauge Theory Gravity (Fact \claimref{synthesis.imported-gtg-dictionary}(i) and (ii)). \emph{Normal form.} Both equations, the Einstein-Cartan equations of \claimref{dark-energy.imported-macdowell-mansouri-theorem-stelle-west}, are built from $\epsilon_{abcd}$, $e$, $\omega$ and $d$, without the metric, which is the polynomiality of \claimref{dark-energy.frozen-point-makes-stelle-west-vector}(1). Under $GL(4)$ on the coordinates and the Lorentz group on the frame, $\epsilon_{abcd}$ is invariant up to $\det$, so ranks and symbol ranks are invariant; a corank-one $e$ may be brought to $e^a = dx^a$ ($a = 0,1,2$), $e^3 = 0$ by a coordinate change and a Lorentz rotation whenever the image is Lorentzian, to $e^0 = 0$ when it is Riemannian, and to $e^0 = e^1 = dx^0$, $e^2 = dx^2$, $e^3 = dx^3$ when it is null; the script runs all three, with identical counts, and the argument below is written for the first. Throughout, a three-form in four variables has four components, one for each omitted coordinate.

\emph{(a) Torsion, pairs not containing $3$.} For $(a,b) \in \{01, 02, 12\}$ the only nonzero terms of $\epsilon_{abcd}\,e^c\wedge T^d$ have $d = 3$, since $e^3 = 0$ removes the terms with $c = 3$; with $c$ the remaining index the equation reads $dx^c\wedge T^3 = 0$. So $T^3\wedge dx^0 = T^3\wedge dx^1 = T^3\wedge dx^2 = 0$. A two-form divisible by $dx^0$ and by $dx^1$ is a multiple of $dx^0\wedge dx^1$, and divisibility by $dx^2$ then forces it to vanish: $T^3 = 0$, six independent conditions. Each of the three equations has one component that is identically zero, the one omitting $c$: three vacated components. \emph{Pairs containing $3$.} For $(a, 3)$ with $\{b, c\} = \{0,1,2\}\setminus\{a\}$ the equation reads $dx^b\wedge T^c = dx^c\wedge T^b$. Split $T^b = \theta^b\wedge dx^3 + \tau^b$ with $\theta^b$ a one-form and $\tau^b$ a two-form in $dx^0, dx^1, dx^2$. The $dx^3$ parts give $dx^b\wedge\theta^c = dx^c\wedge\theta^b$ for all three pairs: writing $\theta^b = \theta^b_\mu dx^\mu$ and $\mu$ for the index other than $b, c$, the components are $\theta^c_\mu = 0$, $\theta^b_\mu = 0$ and $\theta^b_b + \theta^c_c = 0$; over the three pairs every off-diagonal component vanishes and $\theta^0_0 + \theta^1_1 = \theta^0_0 + \theta^2_2 = \theta^1_1 + \theta^2_2 = 0$, so $\theta = 0$: nine conditions, every radical-leg component gone. The $dx^3$-free parts give $dx^b\wedge\tau^c = dx^c\wedge\tau^b$, one condition per pair, three in all, on the nine components of $\tau$: six free, all tangential. Total rank $6 + 9 + 3 = 18$; on the invertible region the same argument with all four $dx^a$ present is the classical one and gives $T = 0$, rank $24$. \emph{Frame equation.} With $e^3 = 0$,
\[
\begin{aligned}
E_3 &= \epsilon_{3bcd}\,dx^b\wedge F^{cd} \qquad (b,c,d \in \{0,1,2\}),\\
E_a &= 2\sum_{b,c}\epsilon_{abc3}\,dx^b\wedge F^{c3} \qquad (a \in \{0,1,2\};\ b, c \in \{0,1,2\}\setminus\{a\}).
\end{aligned}
\]
The unknowns separate: $F^{01}, F^{02}, F^{12}$ enter only $E_3$, and $F^{03}, F^{13}, F^{23}$ enter only $E_0, E_1, E_2$. For $E_3$, the four choices $F^{12} = dx^1\wedge dx^2$, $dx^1\wedge dx^3$, $dx^2\wedge dx^3$ and $F^{02} = dx^2\wedge dx^3$ produce, one each, the four components of a three-form: rank $4$. For $E_0, E_1, E_2$ split $F^{c3} = \phi^c\wedge dx^3 + \psi^c$. The $dx^3$-leg components are $\epsilon_{abc}\,dx^b\wedge\phi^c$, nine from nine, and this map is the $\theta$ system above with the roles renamed, hence bijective; the $dx^3$-free components, one per $E_a$, are hit independently by $\psi^2_{02}$, $\psi^0_{01}$ and $\psi^1_{12}$. Rank $4 + 9 + 3 = 16$, the full target; on the invertible region rank can only be at least as large and the target is sixteen-dimensional. For the radical components, $\iota_r E_a = -\epsilon_{abcd}\,e^b\wedge\iota_r F^{cd}$ is not zero for generic $F$; the script exhibits an integer instance.

\emph{(b) Connection.} An inhomogeneous $de^a$ enters the connection equation only as a source term, which leaves the rank of the linear map $\omega \mapsto \epsilon\,e\wedge T(\omega)$ unchanged, so take $de = 0$: $T^a = \omega^{a}{}_{b}\wedge e^b$. The dying direction: $T^3 = \sum_{b<3}\omega^{3}{}_{b\mu}\,dx^\mu\wedge dx^b$, whose tangential components are the antisymmetric part of $\omega^{3}{}_{b\mu}$ in $(b,\mu)$ and whose radial-leg components are $\omega^{3}{}_{b3}$; $T^3 = 0$ therefore forces the antisymmetric part of $K$ and the three components $\omega^{a3}{}_{3}$ to vanish and leaves the symmetric part of $K$ free, six components; $K$ never enters $T^a$ for $a < 3$, because there it multiplies $e^3 = 0$. The tangential connection: for $a < 3$, $T^a = \omega^{a}{}_{b}\wedge dx^b$ over $b < 3$, so $T^a_{\,3n} = \omega^{a}{}_{n3}$ and the radical-leg condition $\theta = 0$ of (a) forces $\omega^{ab}{}_{3} = 0$, three components; and $T^a_{\,mn} = \omega^{a}{}_{nm} - \omega^{a}{}_{mn}$ for $m,n < 3$ is the classical three-dimensional torsion-contortion map, a bijection of nine-dimensional spaces. Dualising with $\omega^{ab}{}_{\mu} = \epsilon^{abc}M_{c\mu}$ and $t^{a}{}_{d} = \tfrac12\epsilon_{d}{}^{mn}T^a_{\,mn}$, the contraction $\epsilon_{dmn}\epsilon^{amc} = \delta^a_d\delta^c_n - \delta^c_d\delta^a_n$ gives $t = (\operatorname{tr}M)\mathbf{1} - M$, so $t$ is symmetric exactly when $M$ is: the three conditions of (a) are the vanishing of the antisymmetric part of $M$, and the free sector is the symmetric part, six components. Forced: three from each of the antisymmetric part of $K$, the radial-leg $\omega^{a3}{}_{3}$, the radial-leg $\omega^{ab}{}_{3}$ and the antisymmetric part of $M$: rank $12$, free $12$, as claimed.

\emph{(c) Propagation.} The derivative terms of the frame equation are $\epsilon_{abcd}\,e^b\wedge d\omega^{cd}$. For $a = 3$ and $b < 3$ the pair $(c,d)$ is the complement of $b$ in $\{0,1,2\}$, so $\omega^{cd}$ is $\pm M_{b\mu}$: the components of $E_3$ with a $dx^3$ leg are $\partial_3(M_{bx} - M_{xb})$ up to sign, the antisymmetric part only, and the component without a $dx^3$ leg is $\epsilon^{bmn}\partial_m M_{bn}$ up to sign, which vanishes on symmetric $M$. So no derivative of symmetric $M$ appears in $E_3$, and none appears in $E_a$ for $a < 3$, where only $\omega^{c3}$, that is $K$ and its forced complement, occurs. For $K$: the $dx^3$-leg components of $E_a$ ($a < 3$) are $\partial_3 K_{c y}$ with $c$ the index other than $a, b$ and $y \neq b$, which over the two choices of $b$ and the allowed $y$ reach all six components of symmetric $K$; the remaining components carry tangential derivatives of $K$. The script confirms the symbol ranks: $0$ on $M$ in every direction, $6$ on $K$ in the radial direction and $3$ in each tangential direction.

\emph{(d) The Hamiltonian component.} With $de = 0$ and the forced sector at its Levi-Civita value zero, $R^{cd} = \omega^{c}{}_{f}\wedge\omega^{fd}$ and the frame equation is quadratic in $(K, M)$. For $a < 3$ every term of $E_a$ carries $\omega^{c3}$ against $\omega^{3d}$ or a tangential $\omega$, and the $dx^3$-free components need three tangential legs from two one-forms and $dx^b$, which is impossible with $b$ fixed by $a$; the script confirms that all twelve components of $E_a$, $a<3$, vanish identically on the free sector. For $E_3$ on $(\partial_0,\partial_1,\partial_2)$: the $f = 3$ terms give $\epsilon_{3bcd}\,dx^b\wedge\omega^{c}{}_{3}\wedge\omega^{3d} = -\epsilon_{bcd}\,\epsilon^{bxy}K^{c}{}_{x}K^{d}{}_{y} = -\bigl(\delta^x_c\delta^y_d - \delta^x_d\delta^y_c\bigr)K^{c}{}_{x}K^{d}{}_{y} = -\sigma_2'(K)$, and the $f < 3$ terms give, after dualising both factors with $\omega^{cd} = \epsilon^{cdg}M_{g\mu}$, the same contraction on $M$ with the opposite sign, $+\sigma_2'(M)$; cross terms vanish because $\omega^{c}{}_{3}\wedge\omega^{3d}$ and $\omega^{c}{}_{f}\wedge\omega^{fd}$ ($f<3$) lie in different terms. The cosmological term $\ell^{-2}\epsilon_{3bcd}\,e^b\wedge e^c\wedge e^d$ contributes $-6\ell^{-2} = -2\Lambda$ on $(\partial_0,\partial_1,\partial_2)$. The script confirms the identity $E_3(\partial_0,\partial_1,\partial_2) = \sigma_2'(M) - \sigma_2'(K)$ on random elements of the free sector, the vanishing of the cross term, and the coefficient $-6$.
